\documentclass[11pt]{article}

% ============================================================
% Packages
% ============================================================
\usepackage[margin=1in]{geometry}
\usepackage{amsmath,amssymb,mathtools,amsthm}
\usepackage{enumitem}
\usepackage{hyperref}
\usepackage{xcolor}
\usepackage{microtype}
\usepackage{booktabs}
\usepackage{tikz}
\usepackage{cleveref}

% ============================================================
% Hyperref setup
% ============================================================
\hypersetup{
  colorlinks=true,
  linkcolor=blue!60!black,
  citecolor=blue!60!black,
  urlcolor=blue!60!black
}

% ============================================================
% Macros
% ============================================================
\newcommand{\R}{\mathbb{R}}
\newcommand{\F}{\mathcal{F}}
\newcommand{\Pbb}{\mathbb{P}}
\newcommand{\Qbb}{\mathbb{Q}}
\newcommand{\E}{\mathbb{E}}
\newcommand{\dd}{\,\mathrm{d}}
\newcommand{\1}{\mathbf{1}}

% ============================================================
% Theorem environments
% ============================================================
\theoremstyle{plain}
\newtheorem{theorem}{Theorem}[section]
\newtheorem{proposition}[theorem]{Proposition}

\theoremstyle{definition}
\newtheorem{definition}[theorem]{Definition}
\newtheorem{example}[theorem]{Example}

\theoremstyle{remark}
\newtheorem{remark}[theorem]{Remark}

% ============================================================
% Title
% ============================================================
\title{Option Pricing --- Lecture Notes}
\author{Nathan Sauldubois}
\date{MSFE6233 --- Vanilla Options in Practice}

\begin{document}
\maketitle
\vspace{-1em}

\noindent\textbf{Lecture title.}
\emph{Vanilla options in practice: payoffs, Black--Scholes formulas, and static arbitrage bounds}

\vspace{0.5em}
\hrule
\vspace{1em}

\noindent\textbf{Goals of the lecture.}
This lecture is a practical continuation of the no-arbitrage principles introduced previously.
The main objectives are:
\begin{itemize}[leftmargin=1.4em]
    \item to review the basic vanilla option contracts and the standard market vocabulary;
    \item to recall the Black--Scholes pricing formulas and some useful practical approximations;
    \item to derive the main static no-arbitrage inequalities for option prices;
    \item to understand how arbitrage can be constructed when those inequalities are violated.
\end{itemize}

\tableofcontents
\bigskip

% ============================================================
% ============================================================
\section{Introduction and roadmap}
% ============================================================

\subsection{From general no-arbitrage principles to concrete option pricing}

In the previous lecture we introduced the fundamental principles of
arbitrage-free pricing.  
In particular we discussed:

\begin{itemize}[leftmargin=1.4em]
    \item the \emph{dominance principle}: a contract that always pays more
    cannot cost less;
    \item the \emph{law of one price}: two strategies producing the same payoff
    must have the same price;
    \item the idea that replication determines the arbitrage-free price
    of a derivative.
\end{itemize}

These principles provide the conceptual foundation of modern asset pricing.
However, at this stage they remain somewhat abstract.

The goal of the present lecture is to see how these ideas translate
into concrete statements about \emph{option prices}.
In particular we will study the most common derivative contracts
in financial markets: \emph{vanilla options}.

Vanilla options are simple contracts whose payoff depends only on the
terminal value of an underlying asset.
Despite their simplicity, they play a central role in financial markets:
\begin{itemize}[leftmargin=1.4em]
    \item they are actively traded on exchanges,
    \item they are used to hedge risk,
    \item they are the basic building blocks of more complex derivatives.
\end{itemize}

Understanding their pricing and the constraints imposed by no-arbitrage
is therefore a key step toward more advanced models.

\subsection{Why vanilla options matter in practice}

From a practical perspective, vanilla options are the most liquid
derivative instruments in financial markets.

For a given underlying asset (such as a stock index or an equity),
markets typically quote a large set of options with different:
\begin{itemize}[leftmargin=1.4em]
    \item strikes \(K\),
    \item maturities \(T\).
\end{itemize}

These prices form what practitioners call the
\emph{option surface}.
Ensuring that this surface is consistent with no-arbitrage principles
is a fundamental requirement.

In practice, traders and risk managers constantly verify that option prices satisfy a number of structural properties such as:
\begin{itemize}[leftmargin=1.4em]
    \item positivity,
    \item monotonicity in the strike,
    \item convexity with respect to the strike,
    \item consistency across maturities.
\end{itemize}

Violations of these conditions would immediately lead to
\emph{static arbitrage opportunities}.
Understanding these inequalities is therefore both mathematically
interesting and extremely important in practice.

Another reason vanilla options are central is that they admit a
closed-form pricing formula in the Black--Scholes model.
Even though the Black--Scholes model is not perfect,
its formulas remain the standard benchmark used by practitioners.

\subsection{Main themes of the lecture}

This lecture will therefore focus on the practical structure of vanilla option pricing.

We will first review the definition of calls and puts and introduce the
market vocabulary used to describe option positions.
We will then recall the Black--Scholes formulas and discuss how they are
interpreted in practice.

Finally, we will derive the main \emph{static arbitrage constraints}
that option prices must satisfy and see how arbitrage strategies can be
constructed if these constraints are violated.

More precisely, the lecture is organized around four main themes:

\begin{itemize}[leftmargin=1.4em]
    \item \textbf{Payoff geometry.}  
    Understanding option payoffs graphically is essential for building intuition
    and for constructing arbitrage strategies.

    \item \textbf{Market vocabulary.}  
    Traders use a specific terminology (in-the-money, at-the-money,
    straddles, spreads, etc.) which reflects the economic interpretation
    of option positions.

    \item \textbf{Black--Scholes benchmark formulas.}  
    These formulas provide explicit prices for European options
    and serve as a reference point for real market prices.

    \item \textbf{Static arbitrage bounds and consistency checks.}  
    No-arbitrage imposes strong structural restrictions on option prices.
    We will derive these inequalities and understand the trading
    strategies that arise when they are violated.
\end{itemize}

% ============================================================
% ============================================================
\section{Vanilla options: contracts, payoffs, and market language}
% ============================================================

\subsection{European call and European put}

The simplest derivative contracts traded in financial markets are
\emph{European options}.  
A European option gives its holder the right (but not the obligation)
to trade an asset at a predetermined price at a fixed future date.
Let
\begin{itemize}[leftmargin=1.4em]
    \item \(S_T\) be the price of the underlying asset at maturity \(T\),
    \item \(K\) be the \emph{strike price}.
\end{itemize}

\paragraph{European call.}
A European call option gives the holder the right to \emph{buy}
the asset at price \(K\) at time \(T\).

Its payoff is
\[
(S_T-K)^+ = \max(S_T-K,0).
\]

Intuitively, the holder exercises the option only if the asset price
is higher than the strike.

\paragraph{European put.}
A European put option gives the holder the right to \emph{sell}
the asset at price \(K\) at time \(T\).

Its payoff is
\[
(K-S_T)^+ = \max(K-S_T,0).
\]

The option is exercised only if the market price of the asset
is below the strike.

\begin{remark}
European options can only be exercised at maturity \(T\).
This is in contrast with \emph{American options}, which can be exercised
at any time before maturity.
In this course we focus mainly on European options.
\end{remark}

\subsection{Long and short positions}

An option position can be held in two possible directions.

\paragraph{Long position.}
A trader who \emph{buys} an option is said to be \emph{long} the option.
The payoff corresponds exactly to the option payoff defined above.

\paragraph{Short position.}
A trader who \emph{sells} an option is said to be \emph{short} the option.
The payoff is the negative of the option payoff.

For example:
\[
\text{short call payoff} = -(S_T-K)^+.
\]

Being short an option means the trader has sold the right to another market participant.

\subsection{Payoff diagrams}

Payoff diagrams provide a very intuitive way to understand option
contracts and combinations.

\paragraph{Long call.}
The payoff is
\[
(S_T-K)^+.
\]

If \(S_T\le K\), the option expires worthless.  
If \(S_T>K\), the payoff increases linearly with the asset price.

\begin{center}
\begin{tikzpicture}[scale=0.9]
    % axes
    \draw[->] (0,0) -- (5.5,0) node[right] {$S_T$};
    \draw[->] (0,-0.5) -- (0,3) node[above] {Payoff};
    % strike
    \draw[dashed] (2,0) -- (2,0.2);
    \node[below] at (2,0) {$K$};
    % payoff
    \draw[thick] (0,0) -- (2,0) -- (5,2.4);
\end{tikzpicture}
\end{center}

\paragraph{Short call.}
The payoff is
\[
-(S_T-K)^+.
\]

The seller receives the premium initially but faces potentially large losses if the asset price increases.

\begin{center}
\begin{tikzpicture}[scale=0.9]
    % axes
    \draw[->] (0,0) -- (5.5,0) node[right] {$S_T$};
    \draw[->] (0,-3) -- (0,0.8) node[above] {Payoff};
    % strike
    \draw[dashed] (2,0) -- (2,-0.2);
    \node[above] at (2,0) {$K$};
    % payoff
    \draw[thick] (0,0) -- (2,0) -- (5,-2.4);
\end{tikzpicture}
\end{center}

\paragraph{Long put.}
The payoff is
\[
(K-S_T)^+.
\]

This position profits when the underlying price decreases.

\begin{center}
\begin{tikzpicture}[scale=0.9]
    % axes
    \draw[->] (0,0) -- (5.5,0) node[right] {$S_T$};
    \draw[->] (0,-0.5) -- (0,3) node[above] {Payoff};
    % strike
    \draw[dashed] (2,0) -- (2,0.2);
    \node[below] at (2,0) {$K$};
    % payoff
    \draw[thick] (0,2) -- (2,0) -- (5,0);
\end{tikzpicture}
\end{center}

\paragraph{Short put.}
The payoff is
\[
-(K-S_T)^+.
\]

The seller benefits if the asset price remains above the strike.

\begin{center}
\begin{tikzpicture}[scale=0.9]
    % axes
    \draw[->] (0,0) -- (5.5,0) node[right] {$S_T$};
    \draw[->] (0,-3) -- (0,0.8) node[above] {Payoff};
    % strike
    \draw[dashed] (2,0) -- (2,-0.2);
    \node[above] at (2,0) {$K$};
    % payoff
    \draw[thick] (0,-2) -- (2,0) -- (5,0);
\end{tikzpicture}
\end{center}

\begin{remark}
Drawing payoff diagrams is an extremely useful technique.
Many arbitrage arguments rely only on comparing payoff shapes.
\end{remark}
\subsection{Moneyness}

The relation between the current asset price \(S_t\) and the strike \(K\)
is called the \emph{moneyness} of the option.

\paragraph{In the money (ITM).}
An option is in the money if exercising it immediately would generate a positive payoff.

For a call:
\[
S_t > K.
\]

For a put:
\[
S_t < K.
\]

\paragraph{At the money (ATM).}
An option is at the money if the asset price is approximately equal to the strike:

\[
S_t \approx K.
\]

These options are particularly important in practice because they usually carry the largest time value.

\paragraph{Out of the money (OTM).}
An option is out of the money if exercising it immediately would produce zero payoff.

For a call:
\[
S_t < K.
\]

For a put:
\[
S_t > K.
\]

\subsection{Basic option combinations}

More complex payoffs can be constructed by combining several options.
Such combinations are extremely common in practice.

\paragraph{Straddle.}

A \emph{straddle} consists of buying:
\begin{itemize}[leftmargin=1.4em]
    \item one call with strike \(K\),
    \item one put with strike \(K\).
\end{itemize}

The payoff is
\[
(S_T-K)^+ + (K-S_T)^+ = |S_T-K|.
\]

This strategy benefits from large movements of the underlying asset
in either direction.

\paragraph{Protective put.}

A protective put consists of:
\begin{itemize}[leftmargin=1.4em]
    \item holding the underlying asset,
    \item buying a put option.
\end{itemize}

The payoff is
\[
S_T + (K-S_T)^+.
\]

This strategy provides downside protection while maintaining
exposure to upward price movements.

\paragraph{Covered call.}

A covered call consists of:
\begin{itemize}[leftmargin=1.4em]
    \item holding the underlying asset,
    \item selling a call option.
\end{itemize}

The payoff is
\[
S_T - (S_T-K)^+.
\]

This strategy generates additional income from the option premium,
but limits the upside potential of the position.

% ============================================================
% ============================================================
\section{Put--call parity and first applications}
% ============================================================

\subsection{Statement of put--call parity}

Consider a European call and a European put written on the same underlying asset,
with the same strike \(K\) and the same maturity \(T\).
We denote by
\[
C_t \qquad\text{and}\qquad P_t
\]
their prices at time \(t\le T\).

In the no-dividend case, the \emph{put--call parity} relation states that
\[
\boxed{
\ C_t-P_t = S_t - K e^{-r(T-t)}.
}
\]

Equivalently,
\[
\boxed{
\ C_t + K e^{-r(T-t)} = P_t + S_t.
}
\]

\paragraph{Meaning.}
This identity is a static no-arbitrage relation linking the prices of calls and puts.
It does \emph{not} rely on a specific model such as Black--Scholes:
it follows only from the fact that two portfolios with the same terminal payoff
must have the same price today.

\subsection{Proof by replication}

We compare the following two portfolios at time \(t\):

\medskip
\noindent
\textbf{Portfolio A:}
\begin{itemize}[leftmargin=1.4em]
    \item one long call,
    \item cash amount \(K e^{-r(T-t)}\) invested in the bank account.
\end{itemize}

At maturity \(T\), the bank account grows to \(K\), so the terminal payoff of Portfolio A is
\[
(S_T-K)^+ + K.
\]

\medskip
\noindent
\textbf{Portfolio B:}
\begin{itemize}[leftmargin=1.4em]
    \item one long put,
    \item one share of stock.
\end{itemize}

Its terminal payoff is
\[
(K-S_T)^+ + S_T.
\]

\medskip
We now compare the two payoffs state by state.

\paragraph{Case 1: \(S_T\ge K\).}
Then
\[
(S_T-K)^+ + K = (S_T-K)+K = S_T,
\]
while
\[
(K-S_T)^+ + S_T = 0 + S_T = S_T.
\]

\paragraph{Case 2: \(S_T\le K\).}
Then
\[
(S_T-K)^+ + K = 0 + K = K,
\]
while
\[
(K-S_T)^+ + S_T = (K-S_T)+S_T = K.
\]

In both cases, the two portfolios have the same terminal payoff. Therefore,
by the law of one price,
\[
C_t + K e^{-r(T-t)} = P_t + S_t.
\]
This proves put--call parity.

\subsection{Interpretation and practical use}

Put--call parity is one of the most important static relations in option pricing.

\paragraph{A first interpretation.}
The identity says that a call plus discounted cash is economically equivalent
to a put plus the stock:
\[
\text{call} + \text{bond}
\qquad\Longleftrightarrow\qquad
\text{put} + \text{stock}.
\]

So a European call can be viewed as a synthetic position built from a put,
the stock, and the money-market account.

\paragraph{A second interpretation.}
Put--call parity allows us to recover the price of one vanilla option from the other:
\[
C_t = P_t + S_t - K e^{-r(T-t)},
\]
or equivalently
\[
P_t = C_t - S_t + K e^{-r(T-t)}.
\]

\paragraph{Practical use.}
In real markets, put--call parity is used as a consistency check on option quotes.
If quoted prices violate parity, there is a static arbitrage opportunity.
This is one of the first things traders verify when checking market data.

\subsection{Arbitrage if parity is violated}

Suppose the parity relation does \emph{not} hold.

\paragraph{Case 1: the left-hand side is too large.}
Assume
\[
C_t + K e^{-r(T-t)} > P_t + S_t.
\]
Then Portfolio A is overpriced relative to Portfolio B.

An arbitrage strategy is:
\begin{itemize}[leftmargin=1.4em]
    \item sell Portfolio A,
    \item buy Portfolio B.
\end{itemize}

The initial cashflow is strictly positive:
\[
\big(C_t + K e^{-r(T-t)}\big) - \big(P_t+S_t\big) > 0.
\]
At maturity, both portfolios have the same payoff, so the terminal net payoff is \(0\).
Hence this is an arbitrage.

\paragraph{Case 2: the right-hand side is too large.}
Assume instead
\[
C_t + K e^{-r(T-t)} < P_t + S_t.
\]
Then Portfolio B is overpriced relative to Portfolio A.

An arbitrage strategy is:
\begin{itemize}[leftmargin=1.4em]
    \item sell Portfolio B,
    \item buy Portfolio A.
\end{itemize}

Again, the initial cashflow is strictly positive, and the terminal net payoff is \(0\).

\paragraph{Conclusion.}
Any violation of put--call parity immediately creates a static arbitrage opportunity.

\subsection{Optional extension: parity with dividends}

If the underlying asset pays a continuous proportional dividend yield \(q\),
then holding one share of stock is no longer equivalent to holding its full discounted future value.
The put--call parity relation must be adjusted.

In that case, the correct parity becomes
\[
\boxed{
\ C_t-P_t = S_t e^{-q(T-t)} - K e^{-r(T-t)}.
}
\]

Equivalently,
\[
\boxed{
\ C_t + K e^{-r(T-t)} = P_t + S_t e^{-q(T-t)}.
}
\]

\paragraph{Interpretation.}
The factor \(e^{-q(T-t)}\) reflects the fact that the stock pays dividends between \(t\) and \(T\).
So the present value of holding the stock until maturity must be adjusted by the dividend yield.

\begin{remark}
In practice, it is essential to use the correct version of parity:
\begin{itemize}[leftmargin=1.4em]
    \item no dividends \(\Rightarrow\) \(C_t-P_t=S_t-Ke^{-r(T-t)}\),
    \item continuous dividend yield \(q\) \(\Rightarrow\) \(C_t-P_t=S_t e^{-q(T-t)}-Ke^{-r(T-t)}\).
\end{itemize}
Using the wrong parity formula can create the illusion of arbitrage where there is none.
\end{remark}
% ============================================================
% ============================================================
\section{Black--Scholes formulas as practical pricing benchmarks}
% ============================================================

\subsection{Black--Scholes setting}

We now introduce the Black--Scholes model, which provides explicit
pricing formulas for European options.
Although real markets are more complex, these formulas remain the
standard benchmark used by practitioners.

\paragraph{Risk-neutral dynamics.}

Under the risk-neutral measure \(\Qbb\), the stock price is assumed to follow
the stochastic differential equation
\[
\dd S_t = r S_t\,\dd t + \sigma S_t\,\dd W_t,
\]
where
\begin{itemize}[leftmargin=1.4em]
    \item \(r\) is the constant risk-free interest rate,
    \item \(\sigma>0\) is the volatility of the asset,
    \item \((W_t)_{t\ge0}\) is a Brownian motion under \(\Qbb\).
\end{itemize}

\paragraph{Explicit solution.}

This stochastic differential equation has the explicit solution
\[
S_T
=
S_t\exp\!\Big(
(r-\tfrac12\sigma^2)(T-t)+\sigma(W_T-W_t)
\Big).
\]

\paragraph{Lognormal distribution.}

Since \(W_T-W_t\sim \mathcal{N}(0,T-t)\), we obtain that
\[
\log S_T
\]
is normally distributed. In other words, the stock price \(S_T\) follows
a \emph{lognormal distribution}.

This property is the key reason why the Black--Scholes model admits
closed-form option pricing formulas.

\subsection{Closed-form formula for the call}

Consider a European call option with strike \(K\) and maturity \(T\).
Its payoff is
\[
(S_T-K)^+ .
\]

The Black--Scholes formula gives the time-\(t\) price of the call as
\[
\boxed{
C_t
=
S_t\,N(d_1)
-
K e^{-r(T-t)} N(d_2)
}
\]

where \(N(\cdot)\) denotes the cumulative distribution function of the
standard normal distribution:
\[
N(x)=\Pbb(Z\le x), \qquad Z\sim\mathcal{N}(0,1).
\]

The parameters \(d_1\) and \(d_2\) are defined by
\[
d_1=
\frac{\ln(S_t/K)+(r+\tfrac12\sigma^2)(T-t)}
{\sigma\sqrt{T-t}},
\]

\[
d_2=
\frac{\ln(S_t/K)+(r-\tfrac12\sigma^2)(T-t)}
{\sigma\sqrt{T-t}}.
\]

\subsection{Closed-form formula for the put}

The price of a European put option can be obtained either directly
or using put--call parity.

The Black--Scholes formula for the put is
\[
\boxed{
P_t
=
K e^{-r(T-t)} N(-d_2)
-
S_t N(-d_1)
}
\]

An alternative way to compute the put price is to use
put--call parity:
\[
P_t = C_t - S_t + K e^{-r(T-t)}.
\]

In practice, many systems compute the call price first and obtain the
put price using this identity.

\subsection{The parameters \(d_1\) and \(d_2\)}

The quantities \(d_1\) and \(d_2\) play a central role in the
Black--Scholes formula.

They depend on
\begin{itemize}[leftmargin=1.4em]
    \item the ratio \(S_t/K\),
    \item the volatility \(\sigma\),
    \item the time to maturity \(T-t\),
    \item the interest rate \(r\).
\end{itemize}

Observe that
\[
d_2 = d_1 - \sigma\sqrt{T-t}.
\]

Although \(d_1\) and \(d_2\) arise naturally from the mathematical
derivation of the Black--Scholes formula, they also admit a financial
interpretation.

Roughly speaking,
\[
N(d_2)
\]
can be interpreted as the risk-neutral probability that the option
finishes in the money:
\[
\Pbb^{\Qbb}(S_T>K).
\]

Similarly,
\[
N(d_1)
\]
appears when computing the expected value of the stock in the
risk-neutral pricing formula.

% ============================================================
\subsection{Option Greeks}
% ============================================================

In practice, traders are not only interested in the price of an option,
but also in how this price reacts to changes in the model parameters.

These sensitivities are called the \emph{Greeks}.  
They play a central role in risk management and hedging.

Throughout this subsection we consider a European option with price

\[
V(S,t).
\]

The Greeks measure how this price changes when the underlying variables
are modified.

% ------------------------------------------------------------
\paragraph{Delta}
% ------------------------------------------------------------

The \emph{Delta} measures the sensitivity of the option price with
respect to the underlying asset price.

\[
\Delta = \frac{\partial V}{\partial S}.
\]

\paragraph{Interpretation.}

Delta represents the approximate change in the option price for a small
change in the underlying asset price.

For example, if
\[
\Delta = 0.6,
\]
then a small increase of \(1\) in the stock price increases the option
price by approximately \(0.6\).

\paragraph{Black--Scholes formulas.}

For a European call and put:

\[
\Delta_{\text{call}} = N(d_1),
\]

\[
\Delta_{\text{put}} = N(d_1)-1.
\]

\paragraph{Economic interpretation.}

Delta also represents the number of shares needed to hedge a short
option position in the Black--Scholes model.

% ------------------------------------------------------------
\paragraph{Gamma}
% ------------------------------------------------------------

The \emph{Gamma} measures how Delta changes when the stock price moves.

\[
\Gamma = \frac{\partial^2 V}{\partial S^2}.
\]

\paragraph{Interpretation.}

Gamma captures the curvature of the option price with respect to the
underlying asset.

High Gamma means that the Delta changes rapidly as the stock price moves.

\paragraph{Black--Scholes formula.}

For both calls and puts,

\[
\Gamma
=
\frac{N'(d_1)}{S\sigma\sqrt{T-t}},
\]

where \(N'\) denotes the standard normal density.

\paragraph{Remark.}

Gamma is largest for \emph{at-the-money} options and becomes small for
deep in-the-money or deep out-of-the-money options.

% ------------------------------------------------------------
\paragraph{Vega}
% ------------------------------------------------------------

The \emph{Vega} measures the sensitivity of the option price to the
volatility parameter.

\[
\text{Vega} = \frac{\partial V}{\partial \sigma}.
\]

\paragraph{Interpretation.}

Vega represents the change in the option price for a small change in
volatility.

For example, if Vega is equal to \(0.4\), then increasing volatility by
\(1\%\) increases the option price by approximately \(0.004\).

\paragraph{Black--Scholes formula.}

For both calls and puts,

\[
\text{Vega}
=
S \sqrt{T-t}\, N'(d_1).
\]

\paragraph{Remark.}

Vega is typically largest for at-the-money options with intermediate
maturity.

% ------------------------------------------------------------
\paragraph{Theta}
% ------------------------------------------------------------

The \emph{Theta} measures the sensitivity of the option price with
respect to time.

\[
\Theta = \frac{\partial V}{\partial t}.
\]

\paragraph{Interpretation.}

Theta measures the change in the option price as time passes,
all other parameters being fixed.

For long option positions, Theta is typically negative: as time passes,
the option loses part of its time value.
This phenomenon is often called \emph{time decay}.

\paragraph{Black--Scholes formulas.}

For a European call option,

\[
\Theta_{\text{call}}
=
-
\frac{S\sigma N'(d_1)}{2\sqrt{T-t}}
-
rK e^{-r(T-t)} N(d_2).
\]

For a European put option,

\[
\Theta_{\text{put}}
=
-
\frac{S\sigma N'(d_1)}{2\sqrt{T-t}}
+
rK e^{-r(T-t)} N(-d_2).
\]

\paragraph{Remark.}

The first term in the Theta formula corresponds to the decay of the
option's time value due to the reduction of uncertainty as maturity
approaches.

The second term reflects the effect of interest rates on the discounted
strike.

% ------------------------------------------------------------
\paragraph{Rho}
% ------------------------------------------------------------

The \emph{Rho} measures the sensitivity of the option price with respect
to the interest rate.

\[
\rho = \frac{\partial V}{\partial r}.
\]

\paragraph{Interpretation.}

Rho captures how the option price changes when interest rates change.

\paragraph{Black--Scholes formulas.}

\[
\rho_{\text{call}}
=
K(T-t)e^{-r(T-t)} N(d_2),
\]

\[
\rho_{\text{put}}
=
- K(T-t)e^{-r(T-t)} N(-d_2).
\]

% ------------------------------------------------------------
\paragraph{Summary}
% ------------------------------------------------------------

The main Greeks are summarized below:

\[
\begin{array}{c|c}
\text{Greek} & \text{Sensitivity to} \\ \hline
\Delta & \text{underlying price} \\
\Gamma & \text{curvature with respect to the price} \\
\text{Vega} & \text{volatility} \\
\Theta & \text{time} \\
\rho & \text{interest rate}
\end{array}
\]

These quantities are fundamental tools for traders because they describe
how the value of an option reacts to changes in market conditions.



\subsection{How to read the formula in practice}

The Black--Scholes formula also provides valuable intuition about how
option prices depend on the model parameters.

\paragraph{Dependence on the strike.}

For a fixed maturity, the call price decreases as the strike increases.
Intuitively, the right to buy the asset at a higher price is less valuable.

Conversely, the price of a put increases with the strike.

\paragraph{Dependence on maturity.}

For fixed parameters, options with longer maturity tend to be more
expensive because they provide more time for favorable price movements.

This effect is particularly strong for at-the-money options.

\paragraph{Dependence on volatility.}

Volatility plays a crucial role in option pricing.
Higher volatility increases the probability of large movements in the
underlying asset.

Since option payoffs are convex functions of \(S_T\), greater volatility
generally increases the value of both calls and puts.

\paragraph{Dependence on the interest rate.}

Interest rates also influence option prices.

For call options, higher interest rates tend to increase the option value
because the discounted strike \(K e^{-r(T-t)}\) becomes smaller.

For put options, the effect works in the opposite direction.



\subsection{Implied volatility}

In practice, option prices are rarely quoted directly in financial
markets. Instead, traders typically quote the \emph{implied volatility}
of an option.

\paragraph{Definition.}

Given a market price \(C^{\text{mkt}}\) for a European call,
the \emph{implied volatility} is defined as the value of \(\sigma\)
that solves

\[
C^{\text{BS}}(S_t,K,T,r,\sigma) = C^{\text{mkt}},
\]

where \(C^{\text{BS}}\) denotes the Black--Scholes call price.

In other words, the implied volatility is the volatility parameter that,
when inserted into the Black--Scholes formula, reproduces the observed
market price.

\paragraph{Why this convention is used.}

There are several reasons why traders prefer quoting implied
volatility instead of option prices:

\begin{itemize}[leftmargin=1.4em]
\item option prices depend strongly on the level of the underlying asset,
while implied volatility is more stable;

\item volatility provides a natural measure of uncertainty about the
future evolution of the asset price;

\item quoting volatility makes it easier to compare options with
different strikes and maturities.
\end{itemize}

\paragraph{Implied volatility surface.}

In real markets, implied volatility varies across strikes and maturities.
The function

\[
(K,T) \mapsto \sigma_{\text{imp}}(K,T)
\]

is called the \emph{implied volatility surface}.

Studying and modeling this surface is a central topic in modern
quantitative finance.


\paragraph{Volatility smile and volatility skew.}

In the Black--Scholes model, volatility is assumed to be constant.
Therefore the implied volatility should be the same for all strikes \(K\)
and maturities \(T\).

In real markets this is not the case.
When implied volatility is computed from observed option prices,
it typically varies with the strike.

\medskip

\textbf{Volatility smile.}

For some assets (historically in FX markets), the implied volatility
curve as a function of the strike has a roughly symmetric shape:
it is higher for deep in-the-money and deep out-of-the-money options,
and lower near the at-the-money strike.

This pattern is called the \emph{volatility smile}.

\begin{center}
\begin{tikzpicture}[scale=0.9]
    % axes
    \draw[->] (0,0) -- (6,0) node[right] {Strike \(K\)};
    \draw[->] (0,0) -- (0,3.5) node[above] {Implied vol};
    % ATM mark
    \draw[dashed] (3,0) -- (3,0.25);
    \node[below] at (3,0) {ATM};
    % smile
    \draw[thick,domain=0.7:5.3,smooth,variable=\x]
        plot ({\x},{0.8+0.22*(\x-3)^2});
\end{tikzpicture}
\end{center}

\medskip

\paragraph{Volatility skew.}

In equity markets, the pattern of implied volatility across strikes is
typically asymmetric.

More precisely, implied volatility tends to \emph{increase as the strike
decreases}.  
Equivalently, options with low strikes usually have higher implied
volatility than options close to the current spot price.

In particular, \emph{out-of-the-money put options} (which provide
protection against large downward movements of the underlying asset)
often display higher implied volatility than at-the-money options.

This asymmetric pattern is called the \emph{volatility skew}
(or sometimes the \emph{volatility smirk}).

\paragraph{Economic intuition.}

One important explanation comes from the demand for
\emph{downside protection}.

Investors such as asset managers and pension funds often hold large
equity portfolios and are therefore exposed to large market downturns.
Buying out-of-the-money put options provides insurance against such
crashes.

Because these options provide protection against rare but severe
downside events, they are in high demand.
As a result, their market prices tend to be relatively high.

When these higher prices are translated into Black--Scholes implied
volatilities, they appear as higher implied volatility for low strikes.

\paragraph{Interpretation in terms of risk-neutral distributions.}

The volatility skew can also be interpreted in terms of the
risk-neutral distribution of the future asset price.

Recall that option prices contain information about the distribution of
\(S_T\) under the risk-neutral measure.
In particular, the curvature of call prices with respect to the strike
is directly related to the density of this distribution.

A pronounced volatility skew typically indicates that the market
assigns relatively more probability mass to large downward moves of the
underlying asset than predicted by the lognormal distribution of the
Black--Scholes model.

In other words, the market-implied distribution of returns is often
\emph{negatively skewed}.

\paragraph{Contrast with the Black--Scholes model.}

In the Black--Scholes model, volatility is constant and therefore
implied volatility should be the same for all strikes.

The existence of a volatility skew in real markets shows that this
assumption is unrealistic and that the true distribution of asset
returns deviates significantly from the lognormal model.

Understanding and modeling this skew is a major topic in modern option
pricing.

\begin{center}
\begin{tikzpicture}[scale=0.9]
    % axes
    \draw[->] (0,0) -- (6,0) node[right] {Strike \(K\)};
    \draw[->] (0,0) -- (0,3.5) node[above] {Implied vol};
    % ATM mark
    \draw[dashed] (3,0) -- (3,0.25);
    \node[below] at (3,0) {ATM};
    % skew
    \draw[thick,domain=0.7:5.3,smooth,variable=\x]
        plot ({\x},{2.8-0.42*\x+0.035*(\x-3)^2});
\end{tikzpicture}
\end{center}

\medskip

\textbf{Economic interpretation.}

The volatility skew reflects the fact that large downward movements in
equity markets are perceived as more likely than large upward movements.
Investors therefore demand protection against market crashes,
which increases the price of out-of-the-money put options.

\medskip

\textbf{Implication for modeling.}

The presence of volatility smiles and skews shows that the constant
volatility assumption of the Black--Scholes model is not realistic.

This is one of the main motivations for more advanced models such as
local volatility models, stochastic volatility models, and jump models.

% ============================================================
% ============================================================
\section{Practical approximations and trader intuition}
% ============================================================

The Black--Scholes formulas provide exact prices within the model.
However, in practice traders rarely think directly in terms of the full
formula. Instead they rely on simple approximations and intuition about
how option prices behave in different regimes.

Understanding these approximations is very useful for developing a quick
intuition about option prices.

\subsection{Intrinsic value and time value}

The price of an option can be decomposed into two components.

\paragraph{Intrinsic value.}

The intrinsic value is the immediate exercise value of the option.

For a call:
\[
\text{Intrinsic value} = (S_t-K)^+.
\]

For a put:
\[
\text{Intrinsic value} = (K-S_t)^+.
\]

This is the payoff that would be obtained if the option could be exercised
immediately.

\paragraph{Time value.}

The difference between the option price and the intrinsic value is called
the \emph{time value}.

For a call:
\[
\text{Time value} = C_t - (S_t-K)^+ .
\]

For a put:
\[
\text{Time value} = P_t - (K-S_t)^+ .
\]

The time value reflects the possibility that the asset price may move
favorably before maturity.

\begin{remark}
Options that are close to the money typically have the largest time value.
Deep in-the-money or deep out-of-the-money options have relatively small
time value.
\end{remark}

\subsection{Short-maturity intuition}

When the maturity \(T-t\) becomes very small, the option price approaches
its intrinsic value.

Indeed, when there is very little time remaining, there is little
opportunity for the asset price to move significantly.

Therefore,
\[
C_t \approx (S_t-K)^+
\]
for very short maturities.

Similarly,
\[
P_t \approx (K-S_t)^+ .
\]

This approximation becomes increasingly accurate as the maturity
approaches zero.

\subsection{Deep ITM / deep OTM regimes}

The Black--Scholes formula also simplifies in extreme regimes.

\paragraph{Deep in-the-money call.}

If \(S_t \gg K\), the call option behaves approximately like the stock.
Indeed,
\[
C_t \approx S_t - K e^{-r(T-t)}.
\]

In this situation the probability of finishing in the money is close to
one, so the option behaves almost like holding the stock minus the
discounted strike.

\paragraph{Deep out-of-the-money call.}

If \(S_t \ll K\), the call option is very unlikely to finish in the money,
so its value is small.

In that regime
\[
C_t \approx 0.
\]

The option mainly reflects the small probability of a large upward move.

\paragraph{Similar intuition for puts.}

For puts, the opposite regimes hold:
\begin{itemize}[leftmargin=1.4em]
    \item if \(S_t \ll K\), the put is deep in the money and behaves
    approximately like \(K e^{-r(T-t)}-S_t\);
    \item if \(S_t \gg K\), the put price is close to zero.
\end{itemize}

\subsection{ATM options and time value}

At-the-money options (\(S_t \approx K\)) behave differently.

In this regime:
\begin{itemize}[leftmargin=1.4em]
\item the intrinsic value is close to zero,
\item the entire price of the option consists essentially of time value.
\end{itemize}

This explains why at-the-money options are usually the most sensitive to
volatility.

In fact, the Black--Scholes price of an ATM option is roughly of the order
\[
S_t \sigma \sqrt{T-t}.
\]

This scaling reflects the typical magnitude of Brownian fluctuations over
the time horizon \(T-t\).

\subsection{Numerical and practical remarks}

In real markets, option pricing uses a number of standard conventions.

\paragraph{Annualized volatility.}

Volatility is quoted in annualized units.
For example,
\[
\sigma = 20\%
\]
means that the standard deviation of yearly log-returns is approximately
\(20\%\).

\paragraph{Time measured in years.}

Maturity is typically expressed in years.
For example:
\begin{itemize}[leftmargin=1.4em]
\item 3 months \( \approx 0.25 \),
\item 6 months \( \approx 0.5 \),
\item 1 year \( = 1 \).
\end{itemize}

\paragraph{Role of interest rates and dividends.}

Interest rates influence option prices through the discounted strike
\(K e^{-r(T-t)}\).

Dividends reduce the expected future stock price and therefore tend to
decrease call prices while increasing put prices.

\paragraph{Implied volatility.}

In practice, traders rarely quote option prices directly.
Instead they quote the \emph{implied volatility}.

The implied volatility is the value of \(\sigma\) that, when inserted into
the Black--Scholes formula, reproduces the observed market price.

This convention makes it easier to compare options across different
strikes and maturities.
% ============================================================
% ============================================================
\section{Static no-arbitrage inequalities for calls and puts}
% ============================================================

Even without specifying a particular model for the underlying asset,
option prices must satisfy a number of structural inequalities.
These inequalities follow only from the absence of arbitrage.

They provide very useful \emph{consistency checks} for option prices.

Throughout this section we consider European options with strike \(K\)
and maturity \(T\), and denote by \(C_t\) and \(P_t\) the prices of the
call and the put at time \(t\).

% ------------------------------------------------------------
\subsection{Basic upper and lower bounds for calls}
% ------------------------------------------------------------

We begin with simple bounds for the price of a European call.

\paragraph{Lower bound.}

Since the payoff of a call is non-negative,
\[
(S_T-K)^+ \ge 0,
\]
the price of the option must also be non-negative:
\[
C_t \ge 0.
\]

A stronger lower bound can be obtained using put--call parity.

From
\[
C_t - P_t = S_t - K e^{-r(T-t)},
\]
and the fact that \(P_t \ge 0\), we obtain
\[
\boxed{
C_t \ge S_t - K e^{-r(T-t)}.
}
\]

Combining both inequalities gives
\[
\boxed{
C_t \ge \max\!\big(S_t-K e^{-r(T-t)},\,0\big).
}
\]

\paragraph{Upper bound.}

A call option gives the right to buy the asset at price \(K\).
Clearly this right cannot be worth more than the asset itself.

Indeed, if
\[
C_t > S_t,
\]
one could sell the call and buy the stock, obtaining a positive
cashflow initially while never losing money at maturity.

Therefore
\[
\boxed{
C_t \le S_t.
}
\]

\paragraph{Summary.}

We obtain the classical bounds
\[
\boxed{
\max(S_t-K e^{-r(T-t)},0)
\;\le\;
C_t
\;\le\;
S_t.
}
\]

% ------------------------------------------------------------
\subsection{Basic upper and lower bounds for puts}
% ------------------------------------------------------------

Similar arguments yield bounds for the price of a put option.

\paragraph{Lower bound.}

The payoff of a put satisfies
\[
(K-S_T)^+ \ge 0,
\]
so
\[
P_t \ge 0.
\]

Using put--call parity again,
\[
P_t = C_t - S_t + K e^{-r(T-t)},
\]
and the fact that \(C_t \ge 0\), we obtain

\[
\boxed{
P_t \ge K e^{-r(T-t)} - S_t.
}
\]

Combining both inequalities gives

\[
\boxed{
P_t \ge \max\!\big(K e^{-r(T-t)} - S_t,\,0\big).
}
\]

\paragraph{Upper bound.}

A put option allows the holder to sell the asset at price \(K\).
The maximum possible payoff occurs when the asset price is zero.

In that case the payoff is \(K\), so the present value of the maximum
possible payoff is

\[
K e^{-r(T-t)}.
\]

Therefore
\[
\boxed{
P_t \le K e^{-r(T-t)}.
}
\]

\paragraph{Summary.}

We obtain

\[
\boxed{
\max(K e^{-r(T-t)}-S_t,0)
\;\le\;
P_t
\;\le\;
K e^{-r(T-t)}.
}
\]

% ------------------------------------------------------------
\subsection{Positivity of option prices}
% ------------------------------------------------------------

One of the simplest but most fundamental consequences of no-arbitrage
is that option prices must be non-negative.

Indeed, if an option price were negative, a trader could simply
buy the option and receive money immediately.

Since the payoff of both calls and puts is always non-negative,
this would produce an arbitrage.

Therefore

\[
\boxed{
C_t \ge 0
\qquad
\text{and}
\qquad
P_t \ge 0.
}
\]

Although this condition may seem obvious, it is the first consistency
check that any option price must satisfy.

% ------------------------------------------------------------
\subsection{Interpretation of the bounds}
% ------------------------------------------------------------

These inequalities have a simple economic interpretation.

\paragraph{Lower bounds.}

The lower bounds correspond to the fact that an option should be worth
at least its \emph{intrinsic value} (up to discounting effects).

If the option price were lower than this bound, one could construct
an arbitrage by combining the option with the underlying asset
and the risk-free bond.

\paragraph{Upper bounds.}

The upper bounds reflect the fact that an option payoff is limited.

A call cannot be worth more than the underlying asset,
and a put cannot be worth more than the discounted strike.

\paragraph{Practical role.}

In practice these inequalities are constantly used to verify
whether market quotes are consistent.

If any of these bounds is violated, a simple static arbitrage
strategy can immediately be constructed.
% % ============================================================
\section{Monotonicity and convexity in the strike}
% ============================================================

In addition to the basic bounds derived previously, option prices must
also satisfy structural properties with respect to the strike.
These properties follow directly from the shape of option payoffs.

Let \(C(K)\) and \(P(K)\) denote the prices of a European call and put
with strike \(K\) and fixed maturity \(T\).

% ------------------------------------------------------------
\subsection{Monotonicity of call prices with respect to strike}
% ------------------------------------------------------------

Call prices must be decreasing functions of the strike.

\begin{proposition}
If \(K_1 < K_2\), then
\[
C(K_1) \ge C(K_2).
\]
\end{proposition}

\paragraph{Proof.}

The payoff of a call with strike \(K_1\) satisfies
\[
(S_T-K_1)^+ \ge (S_T-K_2)^+
\]
for every value of \(S_T\).

In other words, a call with a lower strike always pays at least as much
as a call with a higher strike.

By the dominance principle, the price must therefore satisfy

\[
C(K_1) \ge C(K_2).
\]

\paragraph{Interpretation.}

A call with a lower strike gives the right to buy the asset at a cheaper
price. It is therefore more valuable.

% ------------------------------------------------------------
\subsection{Monotonicity of put prices with respect to strike}
% ------------------------------------------------------------

The opposite relation holds for put options.

\begin{proposition}
If \(K_1 < K_2\), then
\[
P(K_1) \le P(K_2).
\]
\end{proposition}

\paragraph{Proof.}

For all \(S_T\),

\[
(K_1-S_T)^+ \le (K_2-S_T)^+ .
\]

Thus a put with a higher strike always delivers a payoff at least as
large as a put with a lower strike.

By the dominance principle,

\[
P(K_1) \le P(K_2).
\]

\paragraph{Interpretation.}

A put with a higher strike gives the right to sell the asset at a higher
price, which is more valuable.

% ------------------------------------------------------------
\subsection{Convexity of call prices with respect to strike}
% ------------------------------------------------------------

Option prices must also be convex functions of the strike.

\begin{proposition}
The call price \(C(K)\) is a convex function of the strike \(K\).
\end{proposition}

\paragraph{Proof using convexity of the payoff.}
Fix \(K_1<K_2<K_3\) and assume
\[
K_2=\lambda K_1+(1-\lambda)K_3
\qquad\text{for some } \lambda\in(0,1).
\]
For every fixed value of \(S_T\), the function
\[
K\mapsto (S_T-K)^+
\]
is convex. Therefore,
\[
(S_T-K_2)^+
\le
\lambda (S_T-K_1)^+
+
(1-\lambda)(S_T-K_3)^+.
\]
By the dominance principle, the same inequality must hold for prices:
\[
C(K_2)
\le
\lambda C(K_1)+(1-\lambda)C(K_3).
\]
Hence \(K\mapsto C(K)\) is convex.
\qed

\paragraph{Proof using a butterfly portfolio.}

Consider three strikes
\[
K_1 < K_2 < K_3
\]
and suppose that \(K_2\) lies between the two others.

Define the weights
\[
\lambda=\frac{K_3-K_2}{K_3-K_1},
\qquad
1-\lambda=\frac{K_2-K_1}{K_3-K_1}.
\]

Observe that \(0<\lambda<1\) and that
\[
K_2=\lambda K_1+(1-\lambda)K_3.
\]

Now consider the portfolio

\[
\lambda \text{ calls with strike } K_1
+
(1-\lambda) \text{ calls with strike } K_3
-
1 \text{ call with strike } K_2 .
\]

\paragraph{Payoff at maturity.}

The payoff of this portfolio is

\[
\lambda (S_T-K_1)^+
+
(1-\lambda)(S_T-K_3)^+
-
(S_T-K_2)^+ .
\]

We now examine the payoff in the three possible regions.

\medskip

\textbf{Case 1: \(S_T \le K_1\).}

All calls expire worthless, so the payoff is

\[
0.
\]

\medskip

\textbf{Case 2: \(K_1 \le S_T \le K_2\).}

Only the first call is in the money, so the payoff is

\[
\lambda (S_T-K_1) \ge 0 .
\]

\medskip

\textbf{Case 3: \(K_2 \le S_T \le K_3\).}

Two calls are active and the payoff becomes

\[
\lambda (S_T-K_1) - (S_T-K_2).
\]

Using \(K_2=\lambda K_1+(1-\lambda)K_3\), one can verify that this quantity
is non-negative.

\medskip

\textbf{Case 4: \(S_T \ge K_3\).}

All three options are in the money and the payoff becomes

\[
\lambda (S_T-K_1)
+
(1-\lambda)(S_T-K_3)
-
(S_T-K_2).
\]

A direct computation shows that this expression equals \(0\).

\medskip

Thus the payoff of the portfolio is always non-negative.

\paragraph{No-arbitrage argument.}

Since the terminal payoff is non-negative, the initial cost of the
portfolio must also be non-negative.

Therefore

\[
\lambda C(K_1)
+
(1-\lambda)C(K_3)
-
C(K_2)
\ge 0.
\]

Rearranging gives

\[
C(K_2)
\le
\lambda C(K_1)
+
(1-\lambda)C(K_3),
\]

which is exactly the definition of convexity.

\paragraph{Financial interpretation.}

If this inequality were violated, the butterfly portfolio would have
a negative price but a non-negative payoff, producing an arbitrage.

This type of arbitrage is called a \emph{butterfly arbitrage}.
More precisely, convexity implies

\[
C(K_2)
\le
\frac{K_3-K_2}{K_3-K_1}C(K_1)
+
\frac{K_2-K_1}{K_3-K_1}C(K_3).
\]

If this inequality were violated, a static arbitrage would exist.

% ------------------------------------------------------------
\subsection{Convexity of put prices with respect to strike}
% ------------------------------------------------------------

The same argument applies to put options.

\begin{proposition}
The put price \(P(K)\) is also a convex function of \(K\).
\end{proposition}

This property follows directly from the convexity of the payoff
\[
(K-S_T)^+
\]
with respect to \(K\).

\paragraph{Remark.}

Convexity of option prices is a very important property.
It ensures that the option price curve cannot contain local concave
regions.

In continuous models this property is closely related to the existence
of a non-negative probability density for the underlying asset.

%   

% ============================================================
\section{Maturity inequalities under simple assumptions}
% ============================================================

In addition to constraints with respect to the strike,
option prices must also satisfy consistency conditions across
different maturities.

These relations are sometimes called \emph{calendar inequalities}.
They ensure that option prices behave consistently when the maturity
of the contract changes.

Throughout this section we consider European call options with
fixed strike \(K\) but different maturities.

% ------------------------------------------------------------
\subsection{Monotonicity in maturity when rates and dividends are zero}
% ------------------------------------------------------------

We first consider the simplest situation where
\[
r = 0
\qquad\text{and}\qquad
\text{the asset pays no dividends}.
\]

In that case, call prices must be increasing with maturity.

\begin{proposition}
Let \(T_1 < T_2\).
If interest rates and dividends are zero, then
\[
C(K,T_2) \ge C(K,T_1).
\]
\end{proposition}

\paragraph{Intuition.}

An option with longer maturity gives the holder strictly more
flexibility.

Indeed, holding a call with maturity \(T_2\) allows the investor to
benefit from price movements over a longer time period than a call
with maturity \(T_1\).

Therefore the longer option cannot be less valuable.

\paragraph{Payoff argument.}

Consider two call options with the same strike \(K\)
but maturities \(T_1<T_2\).

At time \(T_1\), the payoff of the first option is

\[
(S_{T_1}-K)^+ .
\]

The holder of the longer option still owns a contract
that remains alive until time \(T_2\).
In particular, the payoff at time \(T_2\) is

\[
(S_{T_2}-K)^+ .
\]

Since the second option offers strictly more opportunities for
favorable price movements, it must have a price at least as large.

\paragraph{Economic interpretation.}

The extra time to maturity increases the probability that the option
finishes in the money.
Therefore longer maturities typically correspond to larger option prices.

\paragraph{Remark (calendar spreads).}

Assume that interest rates and dividends are zero, and consider two
European calls with the same strike \(K\) but different maturities
\[
T_1 < T_2.
\]

Suppose that the maturity monotonicity is violated, that is,
\[
C(K,T_2) < C(K,T_1).
\]

Then the longer-maturity option is cheaper than the shorter-maturity one,
even though it gives strictly more flexibility to its holder.

A natural strategy is the following:
\begin{itemize}[leftmargin=1.4em]
    \item buy one call with maturity \(T_2\),
    \item sell one call with maturity \(T_1\).
\end{itemize}

\paragraph{Initial cashflow.}
At time \(t\), the initial value of this portfolio is
\[
C(K,T_2)-C(K,T_1)<0.
\]
So the trader receives the strictly positive amount
\[
C(K,T_1)-C(K,T_2)>0
\]
at inception.

\paragraph{What happens at time \(T_1\)?}
At time \(T_1\), the short call expires and its payoff is
\[
-(S_{T_1}-K)^+.
\]
The long call with maturity \(T_2\) is still alive. Its market value at
time \(T_1\) is the price of a call with residual maturity \(T_2-T_1\),
namely
\[
C_{T_1}(K,T_2).
\]

Hence the value of the portfolio at time \(T_1\) is
\[
C_{T_1}(K,T_2) - (S_{T_1}-K)^+.
\]

Now, under the same no-arbitrage logic as before, a call with remaining
time to maturity must be worth at least its intrinsic value. Therefore,
\[
C_{T_1}(K,T_2)\ge (S_{T_1}-K)^+.
\]
It follows that
\[
C_{T_1}(K,T_2) - (S_{T_1}-K)^+ \ge 0.
\]

So at time \(T_1\), after the short call has been settled, the trader is
left with a portfolio of non-negative value.

\paragraph{Conclusion.}
The strategy produces:
\begin{itemize}[leftmargin=1.4em]
    \item a strictly positive gain at time \(t\),
    \item and a non-negative value at time \(T_1\).
\end{itemize}
This is an arbitrage.

Such strategies are called \emph{calendar spreads}, because they exploit
inconsistencies between option prices across different maturities.

% ------------------------------------------------------------
\subsection{Why maturity effects can become subtler in general}
% ------------------------------------------------------------

The previous monotonicity property relied on the simplifying
assumption that interest rates and dividends are zero.

When interest rates or dividends are present,
the comparison between different maturities becomes more subtle.

\paragraph{Role of interest rates.}

With positive interest rates, the present value of the strike
\(K e^{-r(T-t)}\) decreases as maturity increases.

This effect tends to increase call prices for longer maturities.

\paragraph{Role of dividends.}

Dividends have the opposite effect.
Since dividends reduce the expected future price of the stock,
they tend to decrease call prices for longer maturities.

As a result, the monotonicity of option prices with respect to maturity
may no longer hold in full generality.

\paragraph{Practical implication.}

In real markets, traders must verify more sophisticated consistency
conditions across maturities in order to rule out
\emph{calendar arbitrage}.

These conditions play an important role in the construction of
arbitrage-free volatility surfaces.

% ============================================================
% ============================================================
\section{How to construct arbitrage when bounds are violated}
% ============================================================

The inequalities derived in the previous sections are not merely
theoretical properties.
Whenever one of these conditions is violated, it is possible to construct
a trading strategy that produces an arbitrage.

In this section we illustrate how such strategies can be constructed.

\subsection{Arbitrage through super-replication}

The same arbitrage logic applies more generally to non-liquid claims.

Suppose that a contingent claim \(\xi\) is traded over the counter and is
not itself liquidly replicable in the market.
Assume however that we can construct a self-financing portfolio of liquidly
traded assets whose terminal wealth \(X_T\) satisfies

\[
X_T \ge \xi
\qquad\text{almost surely.}
\]

In other words, the liquid portfolio \emph{super-replicates} the claim.

Let \(X_0\) denote the initial cost of setting up this portfolio.

\paragraph{Suppose now that the non-liquid claim is sold at price \(P\) with}
\[
P > X_0.
\]

Then the claim is overpriced relative to a super-replicating strategy.

\paragraph{Arbitrage strategy.}

The arbitrageur proceeds as follows:
\begin{itemize}[leftmargin=1.4em]
    \item sell the non-liquid claim \(\xi\) at price \(P\),
    \item use \(X_0\) to buy the liquid super-replicating portfolio,
    \item keep the strictly positive difference
    \[
    P-X_0 > 0
    \]
    as an initial profit.
\end{itemize}

\paragraph{Terminal payoff.}

At maturity, the arbitrageur owes the payoff \(\xi\) to the buyer of the
non-liquid contract.

However, the liquid portfolio delivers \(X_T\), and by construction
\[
X_T \ge \xi.
\]

Therefore, after delivering \(\xi\), the arbitrageur is left with
\[
X_T-\xi \ge 0.
\]

So the terminal net payoff is non-negative in every state of the world.

\paragraph{Conclusion.}

The strategy produces:
\begin{itemize}[leftmargin=1.4em]
    \item a strictly positive gain at time \(0\),
    \item and a non-negative terminal payoff.
\end{itemize}

This is an arbitrage.

\paragraph{Interpretation.}

This argument shows that the price of any claim cannot exceed the initial
cost of a super-replicating strategy built from liquid assets.

Equivalently, if \(V_0\) denotes an admissible arbitrage-free price of the
claim, then necessarily
\[
V_0 \le X_0
\]
for every super-replicating portfolio.

Hence the \emph{super-hedging price} provides an upper bound on the
arbitrage-free price of the claim.


\paragraph{Lower replication (sub-replication).}

A symmetric argument provides a lower bound on the arbitrage-free price
of a contingent claim.

Suppose that we can construct a portfolio of liquid assets with terminal
wealth \(Y_T\) satisfying
\[
Y_T \le \xi
\qquad\text{almost surely.}
\]

In this case the portfolio is said to \emph{sub-replicate} the claim.
Let \(Y_0\) be the initial cost of building this portfolio.

If the market price \(P\) of the claim satisfies
\[
P < Y_0,
\]
then the claim is underpriced relative to the sub-replicating strategy.

\paragraph{Arbitrage strategy.}

An arbitrageur can then:
\begin{itemize}[leftmargin=1.4em]
\item buy the claim \(\xi\) at price \(P\),
\item short the liquid portfolio costing \(Y_0\),
\item keep the initial gain \(Y_0 - P > 0\).
\end{itemize}

At maturity the arbitrageur receives \(\xi\) from the claim and must
deliver \(Y_T\) from the short position.
Since \(Y_T \le \xi\), the remaining payoff
\[
\xi - Y_T \ge 0
\]
is non-negative in every state of the world.

\paragraph{Conclusion.}

Therefore the price of any claim must satisfy
\[
Y_0 \le V_0 \le X_0,
\]
where \(Y_0\) is the cost of a sub-replicating portfolio and \(X_0\) the
cost of a super-replicating portfolio.
These two quantities provide universal no-arbitrage bounds for the
price of the claim.


% ------------------------------------------------------------
\subsection{If a call is too expensive}
% ------------------------------------------------------------

Recall the upper bound
\[
C_t \le S_t .
\]

Suppose instead that
\[
C_t > S_t .
\]

Then the call is overpriced relative to the underlying asset.

\paragraph{Arbitrage strategy.}

\begin{itemize}[leftmargin=1.4em]
\item Sell one call option.
\item Buy one share of the underlying asset.
\end{itemize}

\paragraph{Initial cashflow.}

\[
C_t - S_t > 0.
\]

So the strategy generates an immediate positive profit.

\paragraph{Terminal payoff.}

At maturity:

\begin{itemize}[leftmargin=1.4em]
\item If \(S_T>K\), the call is exercised and the stock is delivered.
The net payoff is \(K\).
\item If \(S_T\le K\), the call expires worthless and the trader still
owns the stock worth \(S_T\ge0\).
\end{itemize}

In both cases the payoff is non-negative.
Therefore the strategy produces an arbitrage.

% ------------------------------------------------------------
\subsection{If a call is too cheap}
% ------------------------------------------------------------

Recall the lower bound
\[
C_t \ge S_t - K e^{-r(T-t)} .
\]

Suppose instead that

\[
C_t < S_t - K e^{-r(T-t)} .
\]

Then the call is underpriced.

\paragraph{Arbitrage strategy.}

\begin{itemize}[leftmargin=1.4em]
\item Buy one call.
\item Sell one share of the underlying asset.
\item Invest \(K e^{-r(T-t)}\) in the risk-free asset.
\end{itemize}

\paragraph{Initial cashflow.}

The initial cost is

\[
C_t - S_t + K e^{-r(T-t)} < 0.
\]

Thus the trader receives money initially.

\paragraph{Terminal payoff.}

At maturity the payoff becomes

\[
(S_T-K)^+ - S_T + K .
\]

It is easy to check that this quantity is always non-negative.
Hence the strategy yields an arbitrage.

% ------------------------------------------------------------
\subsection{If a put is too expensive or too cheap}
% ------------------------------------------------------------

The bounds for put options are

\[
\max(K e^{-r(T-t)}-S_t,0)
\le
P_t
\le
K e^{-r(T-t)}.
\]

\paragraph{If the put is too expensive.}

If

\[
P_t > K e^{-r(T-t)},
\]

one can construct an arbitrage by

\begin{itemize}[leftmargin=1.4em]
\item selling the put,
\item investing the premium in the risk-free asset.
\end{itemize}

The payoff of the put is always bounded by \(K\),
so the strategy cannot lose money.

\paragraph{If the put is too cheap.}

If

\[
P_t < K e^{-r(T-t)} - S_t,
\]

an arbitrage can be constructed using

\begin{itemize}[leftmargin=1.4em]
\item a long put,
\item a long position in the stock,
\item borrowing the discounted strike.
\end{itemize}

% ------------------------------------------------------------
\subsection{If put--call parity is violated}
% ------------------------------------------------------------

Recall the put--call parity relation

\[
C_t - P_t = S_t - K e^{-r(T-t)}.
\]

Suppose instead that

\[
C_t - P_t > S_t - K e^{-r(T-t)} .
\]

Then the portfolio

\[
\text{call} + \text{bond}
\]

is overpriced relative to

\[
\text{put} + \text{stock}.
\]

\paragraph{Arbitrage strategy.}

\begin{itemize}[leftmargin=1.4em]
\item Sell one call and sell the bond \(K e^{-r(T-t)}\).
\item Buy one put and one share of stock.
\end{itemize}

The two portfolios have identical payoffs at maturity,
but the first one is more expensive.
Selling the expensive portfolio and buying the cheap one yields an
arbitrage.

% ------------------------------------------------------------
\subsection{If monotonicity in strike is violated}
% ------------------------------------------------------------

Recall that call prices must satisfy

\[
K_1 < K_2
\quad\Rightarrow\quad
C(K_1) \ge C(K_2).
\]

Suppose instead that

\[
C(K_1) < C(K_2).
\]

Then the higher-strike option is more expensive than the lower-strike one.

\paragraph{Arbitrage strategy.}

\begin{itemize}[leftmargin=1.4em]
\item Buy the call with strike \(K_1\).
\item Sell the call with strike \(K_2\).
\end{itemize}

\paragraph{Terminal payoff.}

The payoff is

\[
(S_T-K_1)^+ - (S_T-K_2)^+ .
\]

It is easy to verify that this payoff is always non-negative.

If the initial cost of the portfolio is negative,
this produces an arbitrage.

% ------------------------------------------------------------
\subsection{If convexity in strike is violated}
% ------------------------------------------------------------

Convexity of call prices implies

\[
C(K_2)
\le
\frac{K_3-K_2}{K_3-K_1}C(K_1)
+
\frac{K_2-K_1}{K_3-K_1}C(K_3).
\]

If this inequality is violated,
one can construct a \emph{butterfly spread}.

\paragraph{Butterfly strategy.}

\begin{itemize}[leftmargin=1.4em]
\item Buy one call with strike \(K_1\),
\item Buy one call with strike \(K_3\),
\item Sell two calls with strike \(K_2\).
\end{itemize}

\paragraph{Payoff.}

The payoff is

\[
(S_T-K_1)^+ - 2(S_T-K_2)^+ + (S_T-K_3)^+.
\]

This payoff is always non-negative.

If the initial price of this portfolio is negative,
the strategy yields an arbitrage.

\paragraph{Interpretation.}

Violations of convexity correspond to
\emph{butterfly arbitrage}.
Such arbitrages are one of the main consistency checks
used when constructing option price surfaces.
% \section{Lecture summary}
% ============================================================

\subsection{Main takeaways}
\begin{itemize}[leftmargin=1.4em]
    \item vanilla options are best understood through payoff geometry;
    \item put--call parity is a central static no-arbitrage identity;
    \item Black--Scholes gives explicit benchmark prices and practical intuition;
    \item option prices must satisfy strong static arbitrage constraints;
    \item violations of these constraints immediately suggest trading strategies.
\end{itemize}

\subsection{Checklist: what students should know after this lecture}
\begin{itemize}[leftmargin=1.4em]
    \item define and draw the payoff of calls, puts, and straddles;
    \item explain ITM / ATM / OTM terminology;
    \item state and use put--call parity;
    \item recall the Black--Scholes call and put formulas;
    \item state the basic no-arbitrage bounds for calls and puts;
    \item explain monotonicity and convexity in the strike;
    \item construct simple arbitrages when static bounds are violated.
\end{itemize}

% ============================================================
\section{Lecture summary}
% ============================================================

\subsection{Main takeaways}

In this lecture we introduced the main practical principles governing
the pricing of vanilla options.  
Although many option pricing models exist, several structural properties
hold independently of any specific model.

The key ideas are summarized below.

\begin{itemize}[leftmargin=1.4em]
    \item \textbf{Payoff geometry is fundamental.}  
    Understanding option payoffs graphically is often the simplest way
    to reason about derivatives and construct arbitrage strategies.

    \item \textbf{Put--call parity is a central identity.}  
    It links the prices of calls, puts, the underlying asset and the
    risk-free asset through a static no-arbitrage argument.

    \item \textbf{Black--Scholes provides a benchmark.}  
    Even though real markets are more complex, the Black--Scholes formulas
    give explicit prices and valuable intuition about how option prices
    depend on volatility, maturity, and the strike.

    \item \textbf{Option prices satisfy strong structural constraints.}  
    These include positivity, upper and lower bounds, monotonicity with
    respect to the strike, convexity in the strike, and consistency
    across maturities.

    \item \textbf{Violations of these constraints imply arbitrage.}  
    Whenever one of these inequalities is violated, it is possible to
    construct a static trading strategy that produces a risk-free profit.
\end{itemize}

\subsection{Checklist: what students should know after this lecture}

After this lecture, students should be able to:

\begin{itemize}[leftmargin=1.4em]
    \item define European calls and puts and draw their payoff diagrams;

    \item explain the notions of \emph{in-the-money}, \emph{at-the-money},
    and \emph{out-of-the-money};

    \item describe simple option combinations such as straddles,
    protective puts, and covered calls;

    \item state and apply the put--call parity relation;

    \item recall the Black--Scholes pricing formulas for European calls
    and puts and interpret the parameters \(d_1\) and \(d_2\);

    \item state the basic no-arbitrage bounds for option prices;

    \item explain why call prices decrease with the strike and why
    option prices are convex in the strike;

    \item identify simple arbitrage opportunities when these
    static constraints are violated;

    \item perform basic consistency checks on option prices
    observed in the market.
\end{itemize}

\paragraph{Looking ahead.}

In the next lectures we will build on these principles to study
dynamic hedging strategies and the mathematical foundations of
continuous-time option pricing models.


\appendix
\section{Option Greeks}


\end{document}